\scanpage{43}

% letterhead
U N I V E R S I T E\\
DES SCIENCES ET TECHNIQUES\\
DU LANGUEDOC\\
\rule{3em}{0.4pt}

Montpellier, le 15 Juin 1978\\

MR/FP\\
\uline{Diffusion}: Liste A\\
n\textsuperscript{o} 3762

% printed form
\uline{AVIS DE SOUTENANCE D'UNE THESE DE DOCTORAT}\\
\uline{D'UNIVERSITE}

\uline{Nom et Prénom du candidat} : GOUSSOPOULOS Jean

\uline{Titre de la thèse} : Etude de la croissance relative postnatale chez le poulet. Principaux tissus et organes, muscles et os individuels. Influence de l'hypophysectomie.

\uline{Date, heure et lieu de soutenance prévus} : 5 Juillet 1978 à 9 heures - U.S.T.L. MONTPELLIER - ENSAM. Laboratoire Zootechnie.

\uline{Nom du Directeur de thèse} : L. DAUZIER

\uline{Désignation du laboratoire où la thèse est préparée} : Station de Physiologie Animale - INRA - MONTPELLIER

\uline{Analyse de l'ouvrage} : Etude descriptive par l'utilisation d'allométrie, de la croissance pondérale relative chez le poulet : principaux tissus et organes ; os et muscles individuels. Mise en évidence de gradients de croissance et de différences de développement selon le sexe. Effets de l'hypophysectomie, dont l'influence est surtout nette au niveau du muscle et du tissus gras.

\uline{Durée approximative des recherches ayant abouti à la thèse} : 5 ans

\uline{Composition du jury proposé} : \quad L. DAUZIER, Président\\
\hspace*{12em}C. THIBAULT\\
\hspace*{12em}J.D. BAYLE\\
\hspace*{12em}M. BENEVENT

Le Secrétaire Général,\\

\handwritten{\unsure{E. Siau}}\\

\uline{E. SIAU.}

% archive sheet number: 42
